\section{A finite experiment with pooled reciprocal directions}
\label{sec:pooled-finite}
The mean-exit argument includes continuous angular distributions, but
that exact assertion alone does not give a finite quadrature algorithm.
Here we specify a finite-support law for which every transition is an
exact grid operation. Direction labels need not be retained with the
output, and no infinite-dimensional physical candidate search occurs.

Fix $t=\min(1,d_0/4)$ and the scale bound $\sigma_0$ in
\eqref{eq:design}. Use the four-direction mixture
\begin{equation}\label{eq:compass-law}
 \rho_\square=\tfrac14(\delta_{te_1}+\delta_{-te_1}
                        +\delta_{te_2}+\delta_{-te_2}).
\end{equation}
It has zero mean, second moment $t^2$ and no common positive drift
projection. Set
\begin{equation}\label{eq:pooled-design}
 H_*=(D_0+2\sigma_0+t)^2/t^2,\quad
 L_*=\lceil2H_*\rceil,\quad n_*=6L_*,\quad
 M_\sigma=\lceil32t/\sigma\rceil,\quad h_\sigma=t/M_\sigma.
\end{equation}
The grid $h_\sigma\Z^2$ has covering radius at most $\sigma/8$,
and all four displacements are grid translations by $M_\sigma$ indices.
These choices depend only on the priors and requested scale. In
particular $n_*$ does not grow as $\sigma$ decreases.

\begin{theorem}[Constructive pooled-direction recovery]
\label{thm:pooled-recovery}
Fix all numerical bounds of $\mathcal B_\eta$, including its known
positive patch margin $\eta$. For sufficiently small $\nu>0$ take
$\sigma=c\nu^{3/2}$ and the two reciprocal protocols
\eqref{eq:reciprocal-forcing} with law \eqref{eq:compass-law}.
At each commanded grid center estimate only the two \emph{pooled}
mean bits. No direction-resolved probabilities or collision positions
are observed. Solid-start attempts and free misses are both zero and
both are counted in the denominator.

A fixed-aperture finite algorithm constructs a smooth periodic
presentation with matched $C^2$ support error, period-basis error and
free-area error at most $C\nu$, and identifies the primitive orbit
count and exact rational generator relations with confidence
$1-\delta$. There are $O(\nu^{-3})$ spatial centers and at most
\begin{equation}\label{eq:pooled-budget}
                  C\nu^{-3}\log\frac{C}{\nu\delta}
\end{equation}
attempts. Constants depend on the stated class, including $\eta$.
Each ideal mean may have independently certified implementation bias
at most $1/(512n_*)$. For example, the position/timing/angle coupling
of Theorem~\ref{thm:calibration} satisfies this budget when
$\ell+\tau+t\alpha\le c\sigma/n_*$. A reciprocal-law total-variation
error can instead be charged as in \eqref{eq:operator-calibration}.

After the means are supplied, $O(\nu^{-6})$ arithmetic, comparison and
fixed-kernel evaluations suffice in the model of
Theorem~\ref{thm:constructive}. This includes $O(n_*\nu^{-3})$
local transition updates, not an exponentially growing tree of
randomized paths. No statistical optimality or Turing bit-complexity
claim is part of this theorem. Exact rational relations do not make
the estimated Euclidean basis coordinates exact.
\end{theorem}
\begin{proof}
Let the target grid cover $U=Q_{R_0}$ with $R_0$ as in
\eqref{eq:aperture}. Command both pooled means on a surrounding square
$W$ containing $U+\overline B_{D_0+2\sigma_0+t+1}$.
The actual forward and translated reverse starts lie in a further
collar of width $t+\sigma_0$. This fixed aperture is independent of
$\nu$, and has $N=O(\sigma^{-2})$ grid nodes.

Estimate each actual mean to accuracy $1/(512n_*)$. Together with the
allowed bias its error relative to the ideal mean is at most
$1/(256n_*)$. Their difference $\widehat g$ therefore has error at
most $1/(128n_*)$. Starting from zero, apply the clipped finite-grid
recursion
\[
 \widehat V_{j+1}(x)=\left[
   \tfrac14\sum_{a\in\{\pm te_1,\pm te_2\}}
        \widehat V_j(x+a)-\widehat g(x)\right]_{[0,1]},
 \qquad j<n_*,
\]
using zero extension outside $W$. No wraparound boundary is used.
The four transitions land exactly on the grid. On that grid, the
true values $u_\sigma(x)$ have the same containing-set and mean-exit
bounds as in the continuum. Proposition~\ref{prop:fixed-exit-aperture}
therefore gives error on $U$ at most
\[
       2^{-6}+1/128=3/128.
\]
Another $1/128$ may be reserved for finite numerical updates. The
resulting error $1/32$ is strictly below the threshold allowance $1/8$
in Lemma~\ref{lem:hulls}. Apply that lemma, then
Proposition~\ref{prop:hull-smooth} and Lemma~\ref{lem:locking}.
These provide Hausdorff error $O(\sigma)$, smooth support error
$O(\sigma^{2/3})$, and the claimed discrete and period conclusions.
The same protected component cutoffs, known $\eta$ thresholds and
rational separation bounds are used, not new assumptions of unequal
shapes or absent individual symmetries. The final smoothed free-area
error is $O(\sigma^{2/3})$ as in Theorem~\ref{thm:constructive}.

Each pooled attempt remains a Bernoulli trial even though its input
has been randomized. Hoeffding's inequality and a union bound over the
$2N$ commanded means give the simultaneous accuracy event with
$Cn_*^2\log(CN/\delta)$ attempts per mean. Since $n_*$ is fixed by
the priors, this proves \eqref{eq:pooled-budget}. Independence refers
to fresh attempts, not to all counterfactual outcomes on one segment.

Every Bellman layer uses four lookups per node; only two layers need be
stored. It costs $O(n_*N)$ arithmetic operations. The finite cluster,
hull, support and period stages have the $O(N^2)$ bound proved in
Theorem~\ref{thm:constructive}. Thus the complete arithmetic order is
$O(\sigma^{-4})$. This counting does not conceal a continuous averaging
oracle: the prescribed law has four atoms and the grid is closed under
its shifts before the explicit aperture killing.
\end{proof}

\subsection{Interval implementation and the role of calibration}
The finite transition step admits exact rational enclosures. If
$g(x)\in[g^-(x),g^+(x)]$ on the observed grid, start $V_0^-=V_0^+=0$
and use
\begin{equation}\label{eq:interval-bellman}
 V_{j+1}^-=[T^W V_j^- -g^+]_{[0,1]},\qquad
 V_{j+1}^+=[T^W V_j^+ -g^-]_{[0,1]}.
\end{equation}
Monotonicity implies that the true exact iterate lies between these
bounds. On the protected grid the occupation therefore belongs to
\begin{equation}\label{eq:occupation-enclosure}
 [\,V_n^-,\ \min(1,V_n^++2^{-\lfloor n/L_*\rfloor})\,].
\end{equation}
Empirical proportions, rational confidence allowances and the four
weights give rational endpoints. Thus the local stage can avoid
uncontrolled floating-point rounding. These enclosures are conditional
on the simultaneous mean-accuracy event and the geometric and reciprocal
preparation hypotheses. They do not validate an apparatus or certify
those priors from the bit record. Certified kernel evaluation remains
a separate requirement of the later smooth geometric stage.

For noisy refinement beyond threshold recovery, a valid iteration depth
balances the two terms in \eqref{eq:mean-exit-noise}; depth tending to
infinity at fixed forcing error is not justified. The physical-candidate
bound \eqref{eq:unknown-support-stability} has no such depth factor,
but does not itself make arbitrary measured forcing physically
admissible. This distinction prevents the comparison theorem from being
used as an unjustified stopping certificate for noisy Bellman iterates.

Optimal stopping, killed Green operators and value iteration have a
well-developed theory. The finite-expected-termination notion in
\cite{Bertsekas} concerns stochastic control problems with specified
costs and termination; our unknown is an occupation whose zero set is
not provided. The collision experiment supplies its signed forcing,
and the component bounds supply a termination estimate for the proof.
We use elementary parts of that machinery, not a new general Bellman
theorem. The new observation conclusion is the recovery from two pooled
reciprocal spatial-input functionals without a common directional cone,
and its fixed-aperture, finite-grid realization. The capacity and active
probing comparisons of Section~\ref{sec:hit-comparison} continue to
apply. In particular this is not a uniformly prepared count-law or
passive spectral inverse.
